\section{Sequence models and local error ideals}\label{sec:sequence-models}

Put \(D=Q_\infty=\ell^\infty(Q)/c_0(Q)\), and identify \(Q\)
with its constant sequences. We use square brackets for sequence
classes when the distinction from coordinates matters. The two
purposes of this section are to turn trace identities into uniform
cut comparisons, and to retain local error bounds when approximation
stages and matrix sizes vary.

\subsection{Regularized limits and their model data}

\begin{proposition}[Regularized limits]\label{seq:regularized}
Let \(\omega\) be a free ultrafilter on \(\N\), and let
\(\tau_n\in T(Q)\) be arbitrary. For a bounded positive sequence
\((x_n)\), the formula
\begin{equation}\label{seq:regularization-formula}
 \rho([(x_n)])=
 \sup_{\eps>0}\lim_{n\to\omega}\tau_n((x_n-\eps)_+)
\end{equation}
defines a lower-semicontinuous tracial weight on \(D\).
Its rank functional is
\begin{equation}\label{seq:regularized-rank}
 d_\rho([(x_n)])=
 \sup_{\eps>0}\lim_{n\to\omega}d_{\tau_n}((x_n-\eps)_+).
\end{equation}
If \(\sigma=\lim_{n\to\omega}\tau_n\) in the canonical topology
on \(T(Q)\), then \(\rho|_Q=\sigma\).
\end{proposition}

\begin{proof}
Pull \(\tau_n\) back to \(\ell^\infty(Q)\) by its \(n\)-th
coordinate homomorphism. The compact Hausdorff trace cone has a
unique limit along \(\omega\), say \(\widetilde\rho\).
For every positive \(x\) in the product, the defining topology gives
\[
 \lim_\omega\tau_n((x_n-\eps)_+)
 \leq\widetilde\rho(x),\qquad
 \widetilde\rho((x-\eps)_+)
 \leq\lim_\omega\tau_n((x_n-\eps)_+).
\]
Taking suprema over \(\eps\) proves
\eqref{seq:regularization-formula} on the product.
Every fixed cut of an element of \(c_0(Q)_+\) is eventually zero,
so \(\widetilde\rho\) vanishes on that ideal and descends to a
lower-semicontinuous trace \(\rho\) on the quotient. This also
proves independence of the representative. The same two inequalities
for constant sequences give \(\rho|_Q=\sigma\).

For completeness, rescale the representatives so that
\(0\leq x_n\leq1\), and let the right side of
\eqref{seq:regularized-rank} be \(L\).
For every \(k\) and \(t>0\),
\[
 \tau_n((x_n^{1/k}-t)_+)
 \leq d_{\tau_n}((x_n-t^k)_+).
\]
Regularization and then the supremum over \(k\) show
\(d_\rho(x)\leq L\).
Conversely, choose a continuous function \(g_\eps\) with values in
\([0,1]\), zero on \([0,\eps/2]\), and one on
\([\eps,1]\). For \(0<t<1\),
\[
 \tau_n((g_\eps(x_n)-t)_+)
 \geq(1-t)d_{\tau_n}((x_n-\eps)_+).
\]
The support of \(g_\eps(x)\) is below the support of \(x\).
Thus \(d_\rho(x)\geq\rho(g_\eps(x))\), and regularization followed
by \(t\downarrow0\) gives
\(d_\rho(x)\geq\lim_\omega d_{\tau_n}((x_n-\eps)_+)\).
Taking the supremum over \(\eps\) proves the reverse inequality.
\end{proof}

We call the weights in \Cref{seq:regularized} \emph{regularized
limits}. The coordinate weights need not be bounded, normalized, or
densely finite. Passing to a subsequence is included by taking an
ultrafilter concentrated on its indices.

Let \(Y\) be a point, a half-open interval, or \(\R\), with a
full-support Radon measure \(m\). The measure is finite on a cone
interval and locally finite on \(\R\). Write \(S=S_Y\).
For \(\nu\in T(P)\), the product \(m\otimes\nu\) is the
integrated densely finite trace on
\(C_0(Y)\otimes\Fin(\nu)\), extended to be infinity on positives
outside that ideal. In particular,
\begin{equation}\label{seq:product-rank}
 d_{m\otimes\nu}(f\otimes a)
 =m(\{f>0\})d_\nu(a)
 \qquad(f\in C_0(Y)_+,\ a\in P_+).
\end{equation}
Here a zero factor makes the product zero, including the expression
\(0\cdot\infty\). Full support of \(m\) ensures that a nonzero
positive continuous function has support of positive measure.

\begin{definition}[Model]\label{seq:model}
A model for \((Y,m,F)\) is a homomorphism \(p:S_Y\to D\) such that
\begin{equation}\label{seq:model-equation}
 \rho\circ p=m\otimes F(\sigma),\qquad \sigma=\rho|_Q,
\end{equation}
for every regularized limit \(\rho\).
\end{definition}

We fix a countable collection of positive elementary tensors for
testing this identity. In each basis ideal \(P(V)\), choose a
countable dense positive set, include its positive rational cuts,
and include the less deeply cut elements needed as enlarged tests.
In \(C_0(Y)\), choose countable dense positive sets on a compact
exhaustion and their rational cuts, with enlarged compact supports.
Take all tensor products of these choices, including the zero tensor
and all positive rational cuts of each chosen factor. A cut above a
factor's norm is simply zero, and its rank identity is automatic.
For the line or a varying interval, a compact test and its enlarged support are chosen inside
one common open interval before an exhaustion index is advanced.

\begin{lemma}[Elementary determination]\label{seq:elementary-tests}
A trace \(\theta\in T(S_Y)\) equals \(m\otimes\nu\) if its
ranks on the preceding countable test collection satisfy
\eqref{seq:product-rank}. In particular, checking those ranks for
every regularized limit suffices for \eqref{seq:model-equation}.
\end{lemma}

\begin{proof}
We first extend the rank identities to all elementary tensors, then
identify the ideal where the weights are densely finite, and finally
compare their bounded restrictions inside that ideal.

\emph{Extension of the tests.}
Suppose that positive contractions \(f_j,a_i\) approximate
\(f,a\) within \(\eta\). The perturbation cut inequalities and
their tensor products give
\[
 (f-3\eta)_+\otimes(a-3\eta)_+
 \precsim(f_j-2\eta)_+\otimes(a_i-2\eta)_+
 \precsim f\otimes a.
\]
The middle tensors belong to the test collection when \(\eta\) is
rational; their enlarged versions give the same comparison with any
fixed positive cuts of \(f\) and \(a\) on the left.
Ranks preserve increasing suprema of such cuts. Applying these
comparisons to both \(\theta\) and \(m\otimes\nu\), and letting
\(\eta\downarrow0\), proves \eqref{seq:product-rank} for every
elementary tensor with compact function support. Exhaustion gives
the remaining elementary tensors.

\emph{Identification of the finite ideal.}
Set \(I=\Fin(\nu)\). The finite ideal of \(\theta\) is exactly
\(C_0(Y)\otimes I\). For one inclusion, a compactly supported
function and a compactly Cuntz supported positive element of \(I\)
give a tensor of finite \(\theta\)-rank. Their ideals generate
\(C_0(Y)\otimes I\).
For the other inclusion, let an open rectangle \(V\times U\) be
contained in the support of \(\Fin(\theta)\). Take nonzero
\(f\in C_c(V)_+\) and cut-supported \(a\in P(U)_+\).
Choose \(f_0\in C_c(V)_+\) equal to one on \(\supp(f)\), and
\(b\in P(U)_+\), \(\delta>0\), with
\(a\precsim(b-\delta)_+\). Then
\[
 f\otimes a\precsim(f_0\otimes b-\delta/2)_+.
\]
The larger tensor belongs to \(\Fin(\theta)\), so \(f\otimes a\)
is compactly Cuntz supported there and has finite \(\theta\)-rank by
\eqref{st:finite-rank-cuts}. Formula \eqref{seq:product-rank} and
\(0<m(f>0)<\infty\) imply \(d_\nu(a)<\infty\).
Such elements generate \(P(U)\), so \(P(U)\subseteq I\).
Rectangles exhaust the open support of any ideal, proving the claim.

\emph{Equality on the finite ideal.}
It remains to identify the densely finite traces on this common
ideal; there is nothing to prove if \(I=0\).
Choose a strictly positive \(h\in I\), put \(e_n=(h-1/n)_+\),
and choose increasing positive compactly supported functions
\(f_n\) whose open supports exhaust \(Y\). The hereditary algebras
\[
 H_n=\Her(f_n\otimes e_n)
     =C_0(\{f_n>0\})\otimes\Her(e_n)
\]
increase with dense union in \(C_0(Y)\otimes I\).
Both traces have finite support rank on \(H_n\), because
\(m(f_n>0)d_\nu(e_n)<\infty\). Their restrictions are bounded
and extend normally to \(H_n^{**}\).
For a positive elementary tensor in \(H_n\), two-variable spectral
integration and the already established elementary rank identities
give
\begin{align*}
 \theta(f\otimes a)
 &=\int_0^{\|f\|}\int_0^{\|a\|}
   d_\theta((f-t)_+\otimes(a-u)_+)\,du\,dt\\
 &=\int_0^{\|f\|}\int_0^{\|a\|}
   m(f>t)d_\nu((a-u)_+)\,du\,dt
 =\Bigl(\int_Y f\,dm\Bigr)\nu(a).
\end{align*}
The first integrand is the value of the normal extension on the
product of the two spectral projections. Tonelli's theorem applies
to these nonnegative functions. Boundedness and algebraic tensor
density now give equality of the two traces on \(H_n\).
Choose a positive contractive approximate unit \((v_n)\) of the
common ideal with each \(v_n\) in the union of the \(H_j\).
For every \(x\geq0\) in that ideal, \(v_nxv_n\) belongs to one
of these hereditary algebras, and
\[
 \theta(v_nxv_n)
 =\theta(x^{1/2}v_n^2x^{1/2})\leq\theta(x),
 \qquad v_nxv_n\longrightarrow x.
\]
Lower semicontinuity gives
\(\theta(x)=\lim_n\theta(v_nxv_n)\), and the same argument holds
for \(m\otimes\nu\). Equality on the localizations therefore
gives equality on the common finite ideal.
Both weights are infinity outside it. Thus they agree everywhere.
\end{proof}

\subsection{Uniform comparison and cut sandwiches}

The following elementary estimate is the source of the uniformity
needed when rows of models are combined.

\begin{lemma}[Bounded witnesses into a fixed cut]
\label{seq:uniform-witness}
Let \(a,b\) be positive elements in a matrix algebra over any
\(C^*\)-algebra, and suppose that \(a\precsim(b-\delta)_+\),
where \(\delta>0\). For every \(\eta>0\) there is a comparison
element \(w\) such that
\begin{equation}\label{seq:witness-bound}
 \|w^*bw-a\|<\eta,
 \qquad
 \|w\|^2\leq\frac{\|a\|+\eta}{\delta}.
\end{equation}
The bound is independent of the matrix size.
\end{lemma}

\begin{proof}
Choose \(v\) with
\(\|v^*(b-\delta)_+v-a\|<\eta\), and put
\[
 g_\delta(t)=
 \begin{cases}
  \sqrt{(t-\delta)_+/t},&t>0,\\
  0,&t=0,
 \end{cases}
 \qquad w=g_\delta(b)v.
\]
Then \(w^*bw=v^*(b-\delta)_+v\), while
\(g_\delta(b)^2\leq\delta^{-1}(b-\delta)_+\).
This proves \eqref{seq:witness-bound}.
\end{proof}

\begin{lemma}[Strict comparison from regularized limits]
\label{seq:strict-comparison}
Let \(x,y\in D_+\). If there is \(0<\alpha<1\) such that
\begin{equation}\label{seq:strict-rank-hypothesis}
 d_\rho(x)\leq\alpha d_\rho(y)
 \quad\hbox{for every regularized limit }\rho,
\end{equation}
then \(x\precsim y\). The same assertion holds after finite
matrix amplification and after coordinatewise folding.
\end{lemma}

\begin{proof}
Choose bounded positive representatives \(x_n,y_n\).
If comparison fails, there is \(\eps>0\) with
\((x-2\eps)_+\not\precsim y\).
For every fixed \(\delta>0\), there must be arbitrarily late
coordinates with
\begin{equation}\label{seq:failed-coordinate-comparison}
 (x_n-\eps)_+\not\precsim(y_n-\delta)_+.
\end{equation}
Otherwise \Cref{seq:uniform-witness}, applied for each approximation
accuracy, supplies bounded sequences of witnesses showing
\((x-\eps)_+\precsim y\), a contradiction.

Choose increasing indices \(n_k\) where
\eqref{seq:failed-coordinate-comparison} holds for \(\delta=1/k\).
Pointwise order in \(\Cu(Q)\), from
\eqref{st:pointwise-order}, gives a trace whose rank on the source
is strictly larger than its rank on the target. The smaller rank is
finite. If the larger rank is finite, divide the trace by that rank;
if it is infinite, divide by one plus the smaller rank. Thus we may
choose \(\tau_k\) with
\[
 d_{\tau_k}((x_{n_k}-\eps)_+)\geq1,
 \qquad
 d_{\tau_k}((y_{n_k}-1/k)_+)\leq1.
\]
Take a free ultrafilter on these indices and the corresponding
regularized limit. Formula \eqref{seq:regularized-rank} gives
\(d_\rho(x)\geq1\). For each fixed \(t>0\), the cut
\((y_{n_k}-t)_+\) is eventually below
\((y_{n_k}-1/k)_+\), so the same formula gives
\(d_\rho(y)\leq1\). This contradicts
\eqref{seq:strict-rank-hypothesis}. All comparisons used to select
\(\tau_k\) took place in \(Q\). Matrix folding preserves those
comparisons and ranks by \Cref{st:folding}.
\end{proof}

\begin{proposition}[Sandwich comparisons]\label{seq:sandwich}
Let \(p:S_Y\to D\) be a model, let \(s=f\otimes a\geq0\), and
suppose that \(0<m_0:=m(f>0)<\infty\). For \(r>0\), choose
\(c_r\in Q_+\) representing \(rG[a]\). Then
\begin{equation}\label{seq:uncut-sandwich}
 c_r\precsim p(s)\quad(0<r<m_0),
 \qquad
 p(s)\precsim c_r\quad(r>m_0).
\end{equation}
The following versions have fixed target cuts.

If \(\eta,\delta>0\), \(0<r<m(f>\eta)\), and \(c\) represents
\(rG[(a-\delta)_+]\), then, for every \(0<\kappa<\eta\delta\),
\begin{equation}\label{seq:lower-trimmed-sandwich}
 c\precsim(p(f\otimes a)-\kappa)_+.
\end{equation}
If \(r'>m_0\) and \(\eps>0\), then a representative \(b\) of
\(r'G[a]\) has a positive cut \((b-\kappa)_+\) such that
\begin{equation}\label{seq:upper-trimmed-sandwich}
 (p(f\otimes a)-\eps)_+\precsim(b-\kappa)_+.
\end{equation}
The choices in these two displays can be made independently of
the model \(p\). The complete collection of such comparisons on
the countable compact tests characterizes the model equation.
\end{proposition}

\begin{proof}
For a regularized limit \(\rho\), put \(\sigma=\rho|_Q\) and
\(t=d_{F\sigma}(a)\). The ranks of \(c_r\) and \(p(s)\) are
\(rt\) and \(m_0t\). Their ratio has the required strict scalar
slack in the appropriate direction, so
\Cref{seq:strict-comparison} proves \eqref{seq:uncut-sandwich}.
This reasoning includes \(t=0\) and \(t=\infty\); the separating
normalization in that lemma forces a finite nonzero value whenever
it is needed for a contradiction.

Apply the lower comparison to
\(s'=(f-\eta)_+\otimes(a-\delta)_+\).
Its support lies in the spectral region \(f\otimes a>\eta\delta\),
so elementary functional calculus gives
\(s'\precsim(f\otimes a-\kappa)_+\) whenever
\(0<\kappa<\eta\delta\). Applying \(p\) proves
\eqref{seq:lower-trimmed-sandwich}.

For the upper comparison, choose
\(m_0<r<r'\) and \(\delta>0\) with
\(\|f\|\delta<\eps/2\). The norm perturbation from \(f\otimes a\)
to \(f\otimes(a-\delta)_+\) gives
\[
 (p(f\otimes a)-\eps)_+
 \precsim p(f\otimes(a-\delta)_+)
 \precsim c,
 \qquad [c]=rG[(a-\delta)_+].
\]
Since positive scalar multiplication preserves compact containment,
\[
 [c]\ll rG[a]\leq r'G[a]=[b].
\]
There is therefore \(\kappa>0\) with
\([c]\leq[(b-\kappa)_+]\). This choice uses only the indicated
constant classes, not \(p\).

Conversely, apply any regularized rank to the fixed-cut comparisons.
Let \(\delta,\eta\downarrow0\) and let the rational lower and
upper scalars approach \(m_0\). The resulting inequalities give
\eqref{seq:product-rank} for the pullback trace on every compact
test. Supremum of cuts handles zero or infinite ranks, and compact
exhaustion handles infinite-measure supports.
\Cref{seq:elementary-tests} now gives
\eqref{seq:model-equation}.
\end{proof}

\begin{corollary}[Constant ideals and all traces]\label{seq:all-traces}
For a model \(p\) and every \(s\in S\),
\begin{equation}\label{seq:constant-ideal}
 \Ideal_D(p(s))=\Ideal_D(Q(U_s)).
\end{equation}
Given countably many models \(p_i:S\to D\), there is a separable
subalgebra \(D_0\subseteq D\) containing their ranges and the
constants \(Q\), such that, for every \(\gamma\in T(D_0)\),
\begin{equation}\label{seq:all-trace-model-equation}
 \gamma\circ p_i=m\otimes F(\gamma|_Q)
 \quad\hbox{for every }i.
\end{equation}
It may also contain any prescribed countable subset of \(D\).
\end{corollary}

\begin{proof}
For a nonzero elementary tensor of finite function support, the
upper sandwiches put its image in the ideal generated by the
corresponding constants. The lower trimmed sandwiches recover the
ideals of all cuts of \(a\), hence the whole constant ideal
\(Q(U_a)\).
For general \(s\), the ideal generated by \(s\) in \(S\) is
generated by positive elementary tensors whose rectangle supports
lie in its open support. Their projected supports exhaust \(U_s\).
Applying the elementary lower bounds gives
\(\Ideal_D(Q(U_s))\subseteq\Ideal_D(p(s))\).
Approximating \(s\) by finite sums of elementary tensors from
\(S(U_s)\) and applying the upper bounds gives the reverse inclusion.

For the second assertion, start with the separable algebra generated
by the stated ranges, constants, and extra elements. For every
countable test, every rational cut and scalar in
\Cref{seq:sandwich}, and every model, adjoin the matrix entries of
comparison witnesses at accuracies tending to zero. The resulting
algebra \(D_0\) is separable, and all these comparisons hold in
\(\Cu(D_0)\), not just in \(\Cu(D)\).
For \(\gamma\in T(D_0)\), write \(\sigma=\gamma|_Q\).
The rank of a constant representing \(rG[a]\) is
\(r d_{F\sigma}(a)\). Applying \(d_\gamma\) to the retained
sandwiches therefore gives the elementary rank formula for each
\(\gamma\circ p_i\). The proof of the converse in
\Cref{seq:sandwich} and \Cref{seq:elementary-tests} applies without
requiring \(\gamma\) to be a regularized limit. This proves
\eqref{seq:all-trace-model-equation}.
\end{proof}

\subsection{The ideals of locally small elements}

\begin{definition}[Local error ideal]\label{seq:smallideal}
For an open \(U\subseteq X\), an element \(x=[(x_n)]\in D\)
belongs to \(J(U)\) if, for every \(\eps>0\), there are a positive
element \(c\) compactly Cuntz supported in \(Q(U)\), a sequence
\(r_n>0\) tending to zero, and an index \(N\), such that
\begin{equation}\label{seq:smallideal-bound}
 d_\tau((|x_n|-\eps)_+)
 \leq r_n d_\tau(c)
 \qquad(n\geq N,\ \tau\in T(Q)).
\end{equation}
The controller \(c\), the decay sequence, and \(N\) may depend on
\(x,U,\eps\). Once these have been chosen, the same bound holds
for all traces and all coordinates in the tail.
\end{definition}

The scalar \(r_n\) is strictly positive; consequently the right
side of \eqref{seq:smallideal-bound} is unambiguously infinity
when \(d_\tau(c)=\infty\). No estimate on the uncut element is
part of this definition.

For example, fix a nonzero positive contraction \(a\) compactly
Cuntz supported in \(Q(U)\). In an additional UHF factor choose
projections \(e_n\) of positive rational traces \(t_n\to0\).
After the trace-preserving identification of \Cref{st:folding},
the elements \(x_n=a\otimes e_n\) satisfy
\[
 \|x_n\|=\|a\|,\qquad
 d_\tau((x_n-\eps)_+)
 =t_n d_\tau((a-\eps)_+)\leq t_n d_\tau(a).
\]
Thus \([(x_n)]\) belongs to \(J(U)\) and is nonzero. This is the
kind of smallness that allows a stabilizing map to retain its norm
while the coefficient in its local rank bound tends to zero.

\begin{proposition}[Properties of the local error ideals]
\label{seq:smallideal-properties}
The definition of \(J(U)\) is independent of representatives and
makes it a closed two-sided ideal of \(D\). The assignment is
order preserving, and
\begin{equation}\label{seq:smallideal-support}
 J(U)\subseteq\Ideal_D(Q(U)).
\end{equation}
If \(\rho\) is a regularized limit and its restriction
\(\sigma=\rho|_Q\) is densely finite on \(Q(U)\), then
\begin{equation}\label{seq:smallideal-vanishing}
 \rho|_{J(U)}=0.
\end{equation}
All assertions use the same conventions in finite matrix algebras.
\end{proposition}

\begin{proof}
We recall the cut estimates that will be used. For bounded operators
in a \(C^*\)-algebra and every tracial rank, the following weaker
forms of the spectral cut inequalities suffice:
\begin{align}
 d_\tau((|x+y|-\eps)_+)
 &\leq d_\tau((|x|-\eps/4)_+)
       +d_\tau((|y|-\eps/4)_+),
 \label{seq:sum-cut}\\
 d_\tau((|bx|-\eps)_+),\ d_\tau((|xb|-\eps)_+)
 &\leq d_\tau((|x|-\eps/(2M))_+)
 \quad(\|b\|\leq M, M>0).
 \label{seq:product-cut}
\end{align}
Also \(d_\tau((|x^*|-t)_+)=d_\tau((|x|-t)_+)\), and
\begin{equation}\label{seq:perturbation-cut}
 \|x-y\|<\eps/4
 \quad\Longrightarrow\quad
 d_\tau((|x|-\eps)_+)
 \leq d_\tau((|y|-\eps/2)_+).
\end{equation}
Here is an explicit cut reduction. Set
\[
 x_t=x f_t(|x|),\qquad
 f_t(s)=\begin{cases}(s-t)_+/s,&s>0,\\0,&s=0.\end{cases}
\]
Then \(\|x-x_t\|\leq t\) and
\(|x_t|=(|x|-t)_+\). Use the perturbation comparison
\(\|z-w\|<\eps\Rightarrow(|z|-\eps)_+\precsim|w|\).
For the sum take \(w=x_{\eps/4}+y_{\eps/4}\), and use
\(|w|\precsim|x_{\eps/4}|\oplus|y_{\eps/4}|\).
For a product take \(w=bx_{\eps/(2M)}\), whose absolute value is
Cuntz below \(|x_{\eps/(2M)}|\). For
\eqref{seq:perturbation-cut} take \(w=y_{\eps/2}\), so that
\(\|x-w\|<3\eps/4\). The identities for adjoints give the other
product order. These comparisons follow from the positive cut
perturbation lemma and \([v^*v]=[vv^*]\); applying the additive
rank functional proves the displays for extended weights as well.

If \(x,y\in J(U)\), choose their controllers at tolerance
\(\eps/4\). A folded finite block sum of the controllers is
compactly Cuntz supported in \(Q(U)\). Taking the maximum of the
two decay sequences proves the bound for \(x+y\) by
\eqref{seq:sum-cut}. Equation \eqref{seq:product-cut} proves
closure under multiplication on either side by an arbitrary bounded
sequence. The equality for adjoints proves self-adjointness.
Equation \eqref{seq:perturbation-cut} proves independence of the
representative. For norm closure, first fix \(\eps\), then choose
one element of the approximating sequence within \(\eps/8\) in
quotient norm. Its representatives are eventually within
\(\eps/4\), so its controller at tolerance \(\eps/2\) works for
the desired output cut. Monotonicity in \(U\) follows directly from
the definition.

To prove \eqref{seq:smallideal-support}, fix a cut of \(|x|\),
choose its controller \(c\precsim(b-\delta)_+\) in \(Q(U)\),
and wait until \(r_n\leq1\).
Pointwise order in \(\Cu(Q)\) gives
\((|x_n|-\eps)_+\precsim c\precsim(b-\delta)_+\).
The bounded witnesses of \Cref{seq:uniform-witness} give comparison
with the constant \(b\) in the quotient. Every positive cut of
\(|x|\) is consequently in \(\Ideal_D(Q(U))\), and so is \(x\).

For vanishing, fix \(\eps\) and the same type of controller.
For each coordinate weight,
\begin{equation}\label{seq:controller-trace-bound}
 d_{\tau_n}(c)
 \leq\frac4\delta\,
       \tau_n((b-3\delta/4)_+).
\end{equation}
Indeed, on the spectral region \(b>\delta\), the function
\((b-3\delta/4)_+\) is at least \(\delta/4\).
Put \(h=(b-\delta/2)_+\). Dense finiteness on \(Q(U)\) and
\eqref{st:finite-rank-cuts} give \(\sigma(h)<\infty\), while
the trace-cone convergence inequalities give
\[
 \lim_\omega\tau_n((b-3\delta/4)_+)
 =\lim_\omega\tau_n((h-\delta/4)_+)
 \leq\sigma(h).
\]
The ranks in \eqref{seq:controller-trace-bound} are therefore bounded
on an \(\omega\)-large set. Multiplying by \(r_n\to0\) in
\eqref{seq:smallideal-bound} and using
\eqref{seq:regularized-rank} yields \(d_\rho(|x|)=0\).
Hence \(\rho\) vanishes on \(J(U)_+\).
The proof uses only finite matrix sums and unnormalized ranks, so it
also proves the matrix assertions, with foldings as in
\Cref{st:folding}.
\end{proof}

\subsection{A diagonalization with fixed controllers}

The diagonalizations used below retain two different kinds of control.
A fixed positive controller bounds a local error's rank, while a
fixed positive target cut bounds the norm of a comparison witness.
Both must be chosen before advancing through the approximation
stages. We call each stage's sequence of coordinate approximations
a \emph{row}. Its finite matrix size may depend on the row and on
the coordinate; all matrices are folded by \Cref{st:folding}.

\begin{lemma}[Local diagonalization]\label{seq:diagonal}
Let \(\mathcal A\) be a countable dense
\(\mathbb Q+i\mathbb Q\) \(*\)-subalgebra of \(S=S_Y\),
containing the localized dense sets to be tested.
For each \(j\), let \(L_j:\mathcal A\to D\) be a map, and choose
bounded coordinate representatives \(L_{j,n}(a)\) for every
\(a\in\mathcal A\). Assume that, for each fixed input,
\[
 \limsup_{j\to\infty}\|L_j(a)\|\leq\|a\|,
\]
and that the quotient errors in every fixed rational linearity,
involution, and multiplicativity relation tend to zero as
\(j\to\infty\). Only bounds on each fixed input are required.
Then one can choose \(j(n)\to\infty\) such that
\[
 L(a)=[(L_{j(n),n}(a))]
\]
extends to a contractive homomorphism \(L:S\to D\).

The same diagonal choice can retain the following countable data.
\begin{enumerate}[label=\textup{(\roman*)}]
\item For each \(l\), fix basis opens \(V_l\Subset U_l\) and
row errors \(e_l^{(j)}\in J(V_l)\), for all sufficiently large
\(j\), with chosen representatives \(e_{l,n}^{(j)}\).
Assume \(\sup_j\|e_l^{(j)}\|<\infty\).
Then their diagonal \(e_l=[(e_{l,n}^{(j(n))})]\) belongs to
\(J(U_l)\).
\item For each prescribed comparison, let
\(a^{(j)},b^{(j)}\in D_+\) have uniformly bounded norms and
chosen bounded positive representatives. Suppose
\(a^{(j)}\precsim(b^{(j)}-\delta)_+\), where \(\delta>0\)
is independent of \(j\). If \(a,b\) are the respective
diagonals, then \(a\precsim(b-\delta')_+\) for any
preselected \(0<\delta'<\delta\).
\end{enumerate}

For local errors arising from maps on \(S\), it suffices to
schedule dense localized inputs in these pairs \(V\Subset U\).
The conclusion then extends to every open \(U\) and every input
in \(S(U)\) if the coordinate error maps are uniformly equicontinuous
on bounded sets, or if their diagonals are the values of a continuous
expression in the resulting quotient maps. This includes differences
from fixed equicontinuous coordinate maps and continuous multilinear
errors, with the other variables on countable dense lists.

In particular, a diagonal of models is again a model when their
data agree on the fixed compact tests and their enlarged versions.
This also applies to models on increasing intervals with unchanged
measure on these tests. Every subsequence of a model retains its
model data.
\end{lemma}

\begin{proof}
\emph{The common controller and the coordinate choice.}
Enumerate all algebraic relations and positive rational cut tests.
Reindex the rows so that row \(j\) includes the first \(j\)
requirements and their quotient norm errors tend to zero.
Include the bounds on each prescribed error and comparison element
among these requirements.
Set any as yet unspecified early-row errors to zero.
In the following construction let \(l\) enumerate pairs consisting
of a prescribed local error and a positive rational tolerance
\(\eps_l\); repeat that error's opens and representatives with
each tolerance.
Choose once and for all the controller
\(b_l=b_{V_l,U_l}\) from \Cref{st:controllers}.

In row \(j\), the \(l\)-th error has, after folding, a bound
\[
 d_\tau((|e^{(j)}_{l,n}|-\eps_l)_+)
 \leq r^{(j)}_{l,n}d_\tau(c^{(j)}_l),
 \qquad r^{(j)}_{l,n}\longrightarrow0,
\]
uniformly for \(\tau\in T(Q)\) eventually in \(n\).
The controller \(c^{(j)}_l\) is compactly Cuntz supported in
\(Q(V_l)\). There is therefore a finite integer \(N_{j,l}\) with
\[
 [c^{(j)}_l]\leq N_{j,l}[b_l].
\]
This integer absorbs internal matrix sizes, repeated blocks, and
any other finite multiplicity at that stage. No bound on it as
\(j\to\infty\) is required.

Fix a sequence \(\eta_j\downarrow0\). After the controllers and
the integers \(N_{j,l}\) have been chosen for \(l\leq j\),
choose strictly increasing coordinate thresholds \(n_j\geq j\)
so large that
\begin{equation}\label{seq:diagonal-delay}
 N_{j,l}r^{(j)}_{l,n}\leq\eta_j
 \qquad(n\geq n_j,\ l\leq j).
\end{equation}
At the same threshold impose the first \(j\) norm and algebraic
requirements, with an additional coordinate error at most
\(\eta_j\). This is possible because a quotient norm is the
limsup of the coordinate norms. Choose representatives for the
countable values; linearity of those representatives is unnecessary,
since the linearity relations themselves are among the tests.

For \(n_j\leq n<n_{j+1}\), put \(j(n)=j\) and use row \(j\).
The diagonal norm bounds make each prescribed value a bounded
sequence. For any fixed local
test \(l\), \eqref{seq:diagonal-delay} then gives, eventually,
\[
 d_\tau((|e_{l,n}|-\eps_l)_+)
 \leq\eta_{j(n)}d_\tau(b_l),
 \qquad j(n)\longrightarrow\infty.
\]
The controller on the right is fixed, so this is precisely the
required \(J(U_l)\) bound. Enumerating all rational cut
tolerances proves membership in that ideal. The diagonal norm
requirements give \(\|L(a)\|\leq\|a\|\), and all algebraic
errors vanish. Thus the quotient values extend uniquely to a
contractive homomorphism on \(S\).

\emph{Extension of local errors.}
For an arbitrary open \(U\), the ideals \(S(V)\), with basis
opens \(V\Subset U'\subseteq U\), have dense union in \(S(U)\).
Indeed they generate \(S(U)\) as a closed sum, and finite sums
already lie in one member: replace finitely many pairs
\(V_i\Subset U'_i\) by their finite unions, using the basis
from \Cref{st:controllers}. Thus the union of the scheduled
localized dense sets is dense in \(S(U)\). Each retained error
on those sets lies in \(J(U')\subseteq J(U)\).
Closedness of \(J(U)\) and the assumed continuity extend this
containment to every localized input. For errors with several
variables, approximate each variable on its scheduled dense set,
using the global lists for unrestricted variables.
The continuity used here is either uniform equicontinuity of the
coordinate errors or continuity of their expression in the completed
quotient maps. Continuity of each row separately would not suffice.

\emph{Comparison witnesses with a fixed target cut.}
Suppose a row
comparison has uniformly bounded source \(a^{(j)}\) and uniformly
bounded \(b^{(j)}\), with target \((b^{(j)}-\delta)_+\) and
\(\delta>0\) independent of \(j\). Include both source and target
norms in the coordinate requirements, so their diagonals belong to
the sequence algebra.
For a chosen positive source cut and comparison accuracy, use
\Cref{seq:uniform-witness} with uncut target
\((b^{(j)}-\delta')_+\) and cut size \(\delta-\delta'\).
Its squared norm is bounded by the source norm plus the accuracy,
divided by \(\delta-\delta'\). Lift the witness to a
bounded coordinate sequence with arbitrarily small additional norm
and equation errors, and put its finite-test equation among the
requirements imposed at \(n_j\). This retained target slack gives
bounded witnesses on the diagonal. Enumerating source cuts and
accuracies proves the desired Cuntz comparisons there. Thus the
order of choices is: fixed tests and target cuts, row controllers
and witnesses, comparison multiplicities, and finally the coordinate
thresholds.

For rows of models, the lower bounds
\eqref{seq:lower-trimmed-sandwich} use target thresholds
\(\kappa<\eta\delta\) determined by fixed trimmed input factors.
The upper bounds \eqref{seq:upper-trimmed-sandwich} use cuts of fixed
constant representatives, determined by the scalar slack and the
fixed source tolerance. Both kinds of target threshold are
independent of the row. They can therefore be retained by the
preceding procedure. The converse part of \Cref{seq:sandwich}
identifies the diagonal model data. On increasing intervals the
same argument uses the common compact tests and their unchanged
measure; no condition is imposed on the moving endpoints.
Restriction to any subsequence retains all fixed-cut comparisons,
which proves the final assertion.
\end{proof}

\begin{remark}\label{seq:lifts}
When the resulting quotient map is a homomorphism with nuclear
separable domain, the completely positive lifting theorem provides
\cpc coordinate lifts \cite[Theorem~3.10]{ChoiEffrosLifting}. Those lifts represent the
same quotient values. Consequently all fixed-cut comparisons and
all memberships in \(J(U)\) are unchanged. This permits the
diagonal construction to be carried out first on countable algebraic
data, while later arguments use actual \cpc coordinate maps.
\end{remark}
